% Vector vocabulary for the visual build walkthrough. Page coordinates in cm,
% top-left origin. All illustrations are local TikZ; no external images needed.
\newcommand{\BuildHeader}[4]{%
  \begin{tikzpicture}[remember picture,overlay,x=1cm,y=-1cm]
    \begin{scope}[shift={(current page.north west)}]
      \fill[teal] (0.70,0.34) rectangle (0.84,0.48);
      \node[anchor=north west,inner sep=0pt,text=teal,font=\sffamily\bfseries\fontsize{11}{13}\selectfont]
        at (1.02,0.29) {#1};
      % #2 intentionally unused: no large bold headline below the section name.
      \node[anchor=north west,inner sep=0pt,text width=14.60cm,align=left,text=graytext,font=\sffamily\fontsize{9}{12}\selectfont]
        at (0.70,0.95) {#3};
      \draw[hairline,line width=0.4pt] (0.70,8.58) -- (15.30,8.58);
      \node[anchor=north east,inner sep=0pt,text=graytext,font=\sffamily\fontsize{6.5}{8}\selectfont]
        at (15.30,8.66) {#4};
    \end{scope}
  \end{tikzpicture}}
\newenvironment{BuildCanvas}{%
  \begin{tikzpicture}[remember picture,overlay,x=1cm,y=-1cm]
  \begin{scope}[shift={(current page.north west)}]
}{\end{scope}\end{tikzpicture}}

\newcommand{\VText}[6]{%
  \pgfmathsetlengthmacro{\VTextWidth}{(#3)*1cm}
  \node[anchor=north west,inner sep=0pt,text width=\VTextWidth,align=left,text=#5,font=\sffamily\fontsize{#4}{\fpeval{#4*1.3}}\selectfont]
    at (#1,#2) {\raggedright #6};}
\newcommand{\VLabel}[4]{%
  \node[anchor=north west,inner sep=0pt,text=#3,font=\sffamily\bfseries\fontsize{8}{10}\selectfont]
    at (#1,#2) {#4};}
\newcommand{\VPanel}[5]{%
  \draw[draw=#5!35,fill=#5!4,line width=0.7pt,rounded corners=4pt] (#1,#2) rectangle (#1+#3,#2+#4);}
\newcommand{\VTag}[4]{%
  \node[anchor=center,draw=#3!35,fill=#3!8,rounded corners=3pt,inner xsep=9pt,inner ysep=5pt,text height=7pt,text depth=1pt,text=#3,font=\sffamily\fontsize{9}{11}\selectfont]
    at (#1,#2) {#4};}
\newcommand{\VArrow}[4]{\draw[teal!75,line width=1.1pt,-{Latex[length=2mm]}] (#1,#2) -- (#3,#4);}
% Source/provenance notes remain in the TeX for reference, not on screen.
\newcommand{\BuildSource}[1]{}

% Icons have a 1.4 x 1.4 nominal canvas; drawings follow the page's downwards y.
\newcommand{\VDoc}[4]{%
  \begin{scope}[shift={(#1,#2)},scale=#3]
    \draw[#4,fill=white,line width=1pt,rounded corners=2pt] (0.12,0.05) rectangle (1.15,1.35);
    \fill[#4!12] (0.25,0.25) rectangle (1.02,0.48);
    \foreach \y in {0.67,0.88,1.09} {\draw[#4!55,line width=0.8pt] (0.30,\y) -- (0.96,\y);}
  \end{scope}}
\newcommand{\VRobot}[4]{%
  \begin{scope}[shift={(#1,#2)},scale=#3]
    \draw[#4,line width=1.1pt] (0.70,0.12) -- (0.70,0.34);
    \fill[#4] (0.70,0.08) circle (0.06);
    \draw[#4,fill=#4!7,line width=1pt,rounded corners=3pt] (0.17,0.35) rectangle (1.23,1.12);
    \fill[#4] (0.47,0.66) circle (0.065);
    \fill[#4] (0.94,0.66) circle (0.065);
    \draw[#4,line width=1pt] (0.49,0.92) -- (0.91,0.92);
    \draw[#4,line width=1.1pt] (0.03,0.57) -- (0.03,0.93) (1.37,0.57) -- (1.37,0.93);
  \end{scope}}
\newcommand{\VPerson}[4]{%
  \begin{scope}[shift={(#1,#2)},scale=#3]
    \draw[#4,fill=#4!8,line width=1pt] (0.70,0.37) circle (0.27);
    \draw[#4,fill=#4!8,line width=1pt,rounded corners=3pt] (0.18,1.30) -- (0.18,1.05) .. controls (0.18,0.65) and (1.22,0.65) .. (1.22,1.05) -- (1.22,1.30) -- cycle;
  \end{scope}}
\newcommand{\VLaptop}[4]{%
  \begin{scope}[shift={(#1,#2)},scale=#3]
    \draw[#4,fill=white,line width=1pt,rounded corners=2pt] (0.10,0.10) rectangle (1.30,1.00);
    \fill[#4!9] (0.20,0.23) rectangle (1.20,0.87);
    \draw[#4,line width=1pt] (0,1.16) -- (1.40,1.16) -- (1.30,1.00) (0,1.16) -- (0.10,1.00);
    \draw[#4,line width=0.9pt] (0.36,0.43) -- (0.53,0.55) -- (0.36,0.67) (0.66,0.68) -- (0.98,0.68);
  \end{scope}}
\newcommand{\VDatabase}[4]{%
  \begin{scope}[shift={(#1,#2)},scale=#3]
    \fill[#4!7] (0.12,0.30) rectangle (1.28,1.05);
    \draw[#4,line width=1pt] (0.12,0.30) -- (0.12,1.05) (1.28,0.30) -- (1.28,1.05);
    \draw[#4,fill=white,line width=1pt] (0.70,0.30) ellipse (0.58 and 0.21);
    \draw[#4,line width=1pt] (0.12,0.68) arc[start angle=180,end angle=0,x radius=0.58,y radius=0.21];
    \draw[#4,line width=1pt] (0.12,1.05) arc[start angle=180,end angle=0,x radius=0.58,y radius=0.21];
  \end{scope}}
\newcommand{\VCheck}[3]{\draw[#3,line width=1.8pt,line cap=round,line join=round] (#1-0.17,#2) -- (#1-0.03,#2+0.14) -- (#1+0.24,#2-0.19);}
\newcommand{\VCross}[3]{\draw[#3,line width=1.4pt,line cap=round] (#1-0.13,#2-0.13) -- (#1+0.13,#2+0.13) (#1-0.13,#2+0.13) -- (#1+0.13,#2-0.13);}
\newcommand{\VFolder}[4]{%
  \begin{scope}[shift={(#1,#2)},scale=#3]
    \draw[#4,fill=#4!8,line width=1pt,rounded corners=1pt] (0.08,0.38) -- (0.08,0.20) -- (0.58,0.20) -- (0.72,0.38) -- (1.30,0.38) -- (1.30,1.15) -- (0.08,1.15) -- cycle;
  \end{scope}}
\newcommand{\VBubble}[7]{%
  \VPanel{#1}{#2}{#3}{1.65}{#4}
  \VLabel{#1+0.25}{#2+0.20}{#4}{#5}
  \VText{#1+0.25}{#2+0.67}{#3-0.50}{11}{ink}{#6}
  \VText{#1+0.25}{#2+1.14}{#3-0.50}{8}{graytext}{#7}
}
